\section{Results}
\label{sec:results}

\subsection{Evaluation Protocol}

We calibrate thresholds from 0.25 through 0.60 in increments of 0.01 on
registered PHerc0814 and PHerc1451 patches.
Checkpoint ranking uses calibrated macro-scroll Dice with a minimum per-scroll
gain guard against released M7.
Final qualification evaluates the raw student at five durations
($1{,}024$, $2{,}048$, $3{,}072$, $4{,}096$, and $8{,}192$ samples) and
thresholds 0.38--0.50, with 0.25 retained as a failure control. Its evidence is
16 locked PHerc0139 slices and six blinded PHerc1447 cubes. These locations are
evaluation-only and spatially separated from the PHerc0139 training boxes.
Promotion requires anti-blob behavior, two-sided morphology and coverage, and
human review; no morphology, M7 blend, or line fitter is included in these
model measurements.
An earlier raw v29 candidate failed the PHerc1447 gate; we do not consider it a
viable model.

\subsection{Duration Study}

Table~\ref{tab:duration} reports every committed validation interval through the
completed \observedSamplesText-sample schedule.
All intervals improve both held-out scrolls over released M7.
The response is shallow and non-monotonic; the best is the
\bestSamplesText-sample interval at $\bestMacro$ calibrated macro-scroll Dice.
This behavior justifies retaining the full ladder instead of exposing only the
longest checkpoint.

\begin{table}[t]
\caption{Completed v31 duration ladder. ``Cal.'' is calibrated macro-scroll Dice;
``$t=.40$'' uses a fixed threshold; ``gain'' is calibrated macro gain over the
M7 initial model. All calibrated thresholds equal the 0.25 search boundary.}
\label{tab:duration}
\centering
\small
\begin{tabular}{rrrrr}
\toprule
Samples & Cal. & $t=.40$ & Gain & Min. gain \\
\midrule
\durationRows
\bottomrule
\end{tabular}
\end{table}

The trust-region projection is active for 97.7\% of optimizer updates in the
first interval and 100\% thereafter.
The method therefore reaches its best observed validation score while remaining
on the boundary of its allowed parameter neighborhood.

Every interval selects threshold 0.25, the lower edge of our sweep.
The calibrated scores therefore compare duration, but do not determine the
operating point: the apparent optimum is censored and 0.25 visibly refills the
blobs that the independent gate is designed to catch.

\subsection{Morphology-Gated Selection}

Joint review selects the raw \selectedSamplesText-sample checkpoint at
$T=\selectedThreshold$. At $T=0.42$ it obtains 16/16 exact component-count
matches and already passes the PHerc1447 anti-blob gate. We nevertheless select
0.45 for its cleaner useful-structure/de-blob balance. At that operating point,
15/16 locked slices have exact component counts; the remaining case is a known
scalar ambiguity in which the intended structures are present but one feature
splits while two strokes touch.

On the six PHerc1447 cubes, the selected model has foreground ratio 0.950508
relative to the v15 anti-blob reference, reference recall 0.911324, zero
interior-regression cubes, zero thickness-regression cubes, and a hard
anti-blob pass. Human review finds useful line recovery in difficult views
without restoration of the released-M7 blobs, although several undergrown
views remain explicit failure watches. The full checkpoint hash is in the
machine-readable release recipe.

\paragraph{Independent perceptual audit.}
We also run FLIP~\cite{andersson2020flip} at 67.0206 pixels per degree on the
exact 1,120 displayed candidate/teacher pairs. At $T=0.45$, the 8,192-sample
model improves FLIP on 12/16 locked slices relative to 2,048 samples; total
teacher-only erosion falls from 4,283 to 4,270 pixels and false additions from
1,201 to 1,155. The audit is not unanimous: mean FLIP narrowly prefers 4,096
samples (0.232073 versus 0.233310), while weighted-median pooling narrowly
prefers 8,192 (0.599964 versus 0.600288). We retain this disagreement. FLIP
rewards teacher resemblance, whereas the blind gate and human review also
penalize the teacher-like girth that can weld adjacent wraps.

\begin{figure*}[t]
  \centering
  \figureasset{figure_2_registered_examples}{0.96\linewidth}{2.2in}
  \caption{Reserved registered comparison. Columns will show CT, released M7,
  projected fine-teacher supervision, the selected 8,192-sample output at
  $T=0.45$, and two-sided
  add/remove errors. Rows will cover PHerc0814, PHerc1451, and the blinded
  PHerc1447 gate, with identical crops and thresholds clearly reported.}
  \Description{Reserved registered comparison across held-out and blind scrolls.}
  \label{fig:registered}
\end{figure*}

\subsection{Line-Fitter Evaluation}

In a calibrated $4\!\times\!5\!\times\!5$ study, the line fitter accepted 481
joins and added 12,697 pixels.
Registered high-resolution precision was 92.5\% overall and 98.4\% for joins no
longer than six pixels.
We observed no full-turn fusion and measured a 21\% reduction in atlas overlap.
A larger 21-cube census measured 98.3\% weighted precision.

Adjacent high-resolution connectivity is useful confirmation but is not, by
itself, proof that a join remains on one wrap.
The umbilicus-aware radial gate is our primary cross-wrap safeguard.
The postprocessor evaluation therefore reports false joins by gate, distance, and
radial regime rather than only an aggregate precision.

\subsection{Interpretation}

The duration experiment establishes consistent positive held-out gains under a
saturated M7 trust constraint, while the completed morphology audit establishes
a release candidate. Their different choices are the central result:
overlap-only calibration prefers 7,168 samples and threshold 0.25, but the
de-blob-aware evidence selects 8,192 and 0.45. Table~\ref{tab:duration} is
generated only from committed run artifacts, and the selection record preserves
the independent locked, blind, and perceptual evidence rather than collapsing
them into one score.
